\section{方法边界、工程启示与开放问题}

到这里，三个发现都已有案例：抽象表示跨域复用、多个路径并行计算、未来目标提前约束当前 token。最后一章不再增加“漂亮机制图”，而是检查这些结果能证明什么、不能证明什么，以及它们如何改变系统设计与评估。

\lecturefigure{slide-62-papers-and-methods.jpg}{进一步案例与方法材料：互动论文和 Circuit Tracing 方法论文}{01:08:27}

\subsection{五个不能省略的限制}

第一，replacement model 是近似，reconstruction error 可能承载关键机制。第二，feature label 由人解释，可能不唯一或不稳定。第三，attribution 是 prompt-local 的局部影响，不是固定全局程序。第四，研究展示的是成功案例；互动论文明确讨论方法只在一部分 prompt 上给出有用结果。第五，也是课堂最强调的一点：attention 未被解释，选择、路由与竞争可能正发生在那里。

\begin{warningbox}{不要把 attribution graph 当作“模型源代码”}
它更像实验仪器生成的局部机制草图：包含真实因果线索，也包含替换误差、剪枝、分组与解释者选择。正确使用方式是提出预测、做干预、找反例，并报告无法解释的部分。
\end{warningbox}

\subsection{从机制发现到系统设计}

这些案例并没有直接给出“消灭幻觉”的按钮，却改变了工程问题的表达。与其只调一个拒答阈值，可以分别测量 known-entity 证据、答案证据与 IDK/refusal 竞争；与其相信模型自述，可以用外部 verifier、工具 trace 与反事实干预检查；与其把固定 forward pass 当成唯一计算预算，可以允许 reflection、recurrent loop 或 adaptive compute 在低置信度时增加检查。

系统层还需要把机制信号与产品约束分开。一个内部 feature 即使能预测错误，也可能随模型版本漂移，无法直接成为长期 API；一个干预即使在实验 prompt 上有效，也可能损害其他能力。更稳健的路径是把机制发现转成可重复测试集、监控指标与故障假说，再通过离线评估、红队、回归测试和逐步上线验证。Interpretability 提供诊断分辨率，不替代完整安全工程。

\teachervoice{01:10:05--01:11:33，老师把“如何防止幻觉”称为十亿美元问题，并明确说自己不知道答案。他提出更好 calibration、让 reasoning model 反思、增加自检计算，以及 recurrent/adaptive compute 等方向，同时承认这些选择可能牺牲能力或创造性。}

\begin{importantbox}{工程上最值得带走的四点}
\begin{enumerate}
  \item 把模型解释视为待验证输出，而不是内部事实接口；
  \item 在生成全过程监控不确定性、安全与目标漂移，而非只检查 prompt；
  \item 用反事实 intervention 与回放 trace 验证机制假说；
  \item 对困难 token 分配可变计算，让模型在需要时检索、反思或停止。
\end{enumerate}
\end{importantbox}

\subsection{一套可复用的 circuit-tracing 审计流程}

\begin{lstlisting}[language=Python,caption={从行为问题到可证伪机制假说的审计流程}]
behavior = define_prompt_and_counterfactuals()
features = fit_sparse_replacement_model(behavior.activations)
graph = trace_prompt_specific_attribution(features)
hypotheses = label_prune_and_group(graph)
for hypothesis in hypotheses:
    intervene_in_base_model(hypothesis)
    test_predicted_behavior_change()
report_reconstruction_error_failures_and_scope()
\end{lstlisting}

这段伪代码故意把“report failures”放在最后一个必选步骤。只有成功图而没有失败率，会造成 discovery bias；只有可读 feature 而没有干预，会造成故事偏差；只有单 prompt 干预而没有反事实集合，会把偶然局部效应误当成稳健机制。

实际审计还应保存随机种子、字典版本、剪枝阈值、干预强度和完整 prompt，使另一位研究者能够复现同一张图，并判断结论究竟来自模型机制还是分析超参数。

当不同超参数给出不同路径时，报告这种不稳定性本身就是结果：它说明当前证据只能支持较粗粒度的机制结论，不能安全地解释到单一节点或单一边。

\subsection{开放问题}

课堂问答留下了几条真正困难的研究线：怎样解释 attention 的选择与路由；怎样判断一个 feature 在不同 prompt 和模型中是否“同一个”；怎样量化 graph 对原模型行为的覆盖率；怎样从 token 级机制扩展到长上下文、工具调用与 agent loop；怎样区分必要路径、充分路径与冗余备份路径；怎样让模型在保持创造力的同时更好校准自知之明。

\teachervoice{01:11:33--01:11:57，老师提醒“hallucination”本身并不总是清晰定义：长文本中究竟哪个 token 开始错、错误是事实还是语境偏离，往往没有唯一边界。研究指标必须先定义对象，不能把模糊产品抱怨直接当机制标签。}

\begin{knowledgebox}{研究报告应同时给出三张表}
\textbf{Coverage 表}：哪些 prompt/节点被解释；\textbf{causal 表}：哪些干预按预测改变行为；\textbf{failure 表}：哪些图不稳定、误差大或无法命名。只有三者同时存在，读者才能判断方法的实际适用范围。
\end{knowledgebox}

\subsection{本章小结}

Circuit tracing 把“模型内部可能发生什么”推进到可画图、可干预、可反驳的层次，但它不是完整读取器。当前最有价值的成果，是形成更细的工程问题与实验接口：区分表示、执行与选择，记录误差，验证局部因果，并把不确定性留在结论中。

